% titel: Blatt 1, A5: Differentialgleichungen -- Pendel
% skript: Kap. Differentialgleichungen, harmonischer Oszillator

% Unterschied algebraische Gleichung vs. DGL; Loesung der Oszillator-DGL pruefen.

Eine gewoehnliche Differentialgleichung (DGL) bestimmt eine \emph{Funktion},
keine Zahl; die allgemeine Loesung enthaelt frei waehlbare Konstanten.

\begin{werkzeugkasten}[Werkzeuge]
\textbf{Algebraische Gleichung} (z. B. $x^2+2=0$): gesucht ist eine bzw.
mehrere \emph{Zahlen} $x$, die die Gleichung erfuellen.

\textbf{Differentialgleichung} (z. B. $-D\,x(t)=m\,\ddot x(t)$): gesucht
ist eine \emph{Funktion} $x(t)$, die zusammen mit ihren Ableitungen die
Gleichung identisch erfuellt. Nach $x(t)$ differenziert man mit den Regeln
\begin{equation}
  x=x_0\cos(\omega t+\varphi_0),\quad
  v=\dot x=-x_0\omega\sin(\omega t+\varphi_0),\quad
  a=\ddot x=-x_0\omega^2\cos(\omega t+\varphi_0).
\end{equation}

\textbf{Harmonischer Oszillator:} Jede Funktion der Form
$x(t)=x_0\cos(\omega t+\varphi_0)$ loest die DGL
\begin{equation}
  \ddot x+\omega^2 x=0,
\end{equation}
wenn die Kreisfrequenz $\omega$ passt. Das ist die Grundgleichung der
Schwingungslehre.
\end{werkzeugkasten}

\begin{werkzeugkasten}[Vorgehen]
\begin{itemize}
  \item Ableitungen bilden: $v=\dot x$ als Funktion, nicht Zahl.
  \item Zum Nachweis einsetzen: $\ddot x+\omega^2x$ berechnen und zeigen,
        dass sie null ist.
  \item Die Konstanten $x_0$ (Amplitude) und $\varphi_0$ (Phase) sind die
        Freiheitsgrade; $\omega$ ist durch das System (Feder, Pendel)
        festgelegt.
\end{itemize}
\end{werkzeugkasten}

\begin{achtung}[Typische Fallen]
\begin{itemize}
  \item ``Gesucht'' bei einer DGL ist eine ganze Funktion, nicht ein
        einzelner Wert.
  \item Ableitungen mit der Kettenregel (Argument $\omega t+\varphi_0$),
        nicht einfach $\cos$ differenzieren.
  \item $\varphi_0$ verschiebt nur den Nullpunkt der Zeit, beeinflusst aber
        die Periodendauer nicht.
\end{itemize}
\end{achtung}
